\documentclass[12pt]{article}
\usepackage{../../ficha}

\begin{document}

\ficha{Radicais e Potências de Expoente Racional}

\instrucoes{Resolve sem calculadora. Simplifica os radicais e racionaliza os denominadores. Considera que as letras representam números reais positivos.}

\secao{Simplificação de Radicais}{Simplifica cada radical.}
\begin{exercicios}
  \item $\sqrt{12}$
  \item $\sqrt{50}$
  \item $\sqrt{72}$
  \item $\sqrt{200}$
  \item $\sqrt[3]{16}$
  \item $\sqrt[3]{54}$
  \item $\sqrt[4]{32}$
  \item $\sqrt{18a^3}$
  \item $\sqrt[3]{24x^4}$
  \item $\sqrt{\dfrac{75}{4}}$
\end{exercicios}

\secao{Operações com Radicais}{Calcula e simplifica.}
\begin{exercicios}
  \item $3\sqrt{2} + 5\sqrt{2}$
  \item $\sqrt{12} + \sqrt{27}$
  \item $\sqrt{50} - \sqrt{18} + \sqrt{8}$
  \item $2\sqrt{45} - 3\sqrt{20}$
  \item $\sqrt{3} \times \sqrt{12}$
  \item $\sqrt{6} \times \sqrt{10}$
  \item $\sqrt[3]{4} \times \sqrt[3]{2}$
  \item $\dfrac{\sqrt{60}}{\sqrt{5}}$
  \item $(2\sqrt{3})^2$
  \item $(\sqrt{5} + 1)^2$
  \item $(\sqrt{7} - \sqrt{2})(\sqrt{7} + \sqrt{2})$
  \item $\sqrt{2} \times \sqrt[3]{2}$ \; (escreve como um único radical)
\end{exercicios}

\secao{Racionalização de Denominadores}{Racionaliza os denominadores e simplifica.}
\begin{exercicios}
  \item $\dfrac{1}{\sqrt{2}}$
  \item $\dfrac{6}{\sqrt{3}}$
  \item $\dfrac{5}{2\sqrt{5}}$
  \item $\dfrac{\sqrt{2}}{\sqrt{6}}$
  \item $\dfrac{3}{\sqrt[3]{3}}$
  \item $\dfrac{1}{\sqrt{3} - 1}$
  \item $\dfrac{4}{\sqrt{5} + 1}$
  \item $\dfrac{\sqrt{2}}{\sqrt{2} + 1}$
  \item $\dfrac{\sqrt{3} + 1}{\sqrt{3} - 1}$
  \item $\dfrac{2}{\sqrt{7} - \sqrt{5}}$
\end{exercicios}

\secao{Potências de Expoente Racional}{Escreve na forma de radical ou de potência, conforme pedido, e calcula quando possível.}
\begin{exercicios}
  \item Escreve na forma de radical: $5^{\frac{1}{2}}$, $\;2^{\frac{3}{4}}$, $\;x^{\frac{2}{3}}$.
  \item Escreve na forma de potência: $\sqrt{7}$, $\;\sqrt[3]{a^2}$, $\;\dfrac{1}{\sqrt[5]{x}}$.
  \item $9^{\frac{1}{2}}$
  \item $8^{\frac{2}{3}}$
  \item $16^{\frac{3}{4}}$
  \item $27^{-\frac{1}{3}}$
  \item $\left(\dfrac{4}{9}\right)^{-\frac{3}{2}}$
  \item $32^{0{,}4}$
  \item $2^{\frac{1}{2}} \times 2^{\frac{1}{3}}$
  \item $\dfrac{a^{\frac{3}{4}}}{a^{\frac{1}{4}}}$
  \item $\left(x^{\frac{2}{3}}\right)^{\frac{3}{4}}$
\end{exercicios}

\secao{Equações e Expressões}{Resolve ou simplifica.}
\begin{exercicios}
  \item Resolve: $x^2 = 18$
  \item Resolve: $x^3 = -64$
  \item Resolve: $\sqrt{2}\,x = 6$
  \item Resolve: $x^{\frac{1}{2}} = 5$
  \item Resolve: $x^{\frac{2}{3}} = 4$
  \item Simplifica: $\dfrac{\sqrt{8} + \sqrt{18}}{\sqrt{2}}$
  \item Simplifica: $\sqrt{a\sqrt{a}}$, escrevendo como potência de $a$.
\end{exercicios}

\secao{Desafios}{Resolve os problemas.}
\begin{exercicios}
  \item Um quadrado tem área 48\,cm$^2$. Determina o lado e a diagonal, na forma de radical simplificado.
  \item Mostra que $\dfrac{1}{\sqrt{2} + 1} = \sqrt{2} - 1$.
  \item Calcula $\dfrac{1}{\sqrt{1} + \sqrt{2}} + \dfrac{1}{\sqrt{2} + \sqrt{3}} + \dfrac{1}{\sqrt{3} + \sqrt{4}}$.
  \item Um cubo tem volume 54\,cm$^3$. Determina a aresta e a área total, na forma de radical simplificado.
  \item Ordena por ordem crescente, sem calculadora: $\sqrt{3}$, $\sqrt[3]{5}$, $\sqrt[6]{24}$. (Sugestão: eleva à sexta potência.)
  \item Mostra que $\sqrt{3 + 2\sqrt{2}} = 1 + \sqrt{2}$.
\end{exercicios}

\end{document}
